\begingroup
\clearpage\ArithmeticHead{335}{R}{{\scriptsize ZUR THEORIE DER REIHE $F(\alpha,\beta,\gamma,x)$.\quad II.}}
\NumberedItem{5}
\noindent{\small[S. 246]}\par
\begin{center}\textls[120]{Berechnung von} $\displaystyle\log\varpi\left[=\log_{10}\int_0^1\frac{dx}{\surd(1-x^4)}\right]$\end{center}
\[\begin{aligned}
[\varpi=]\;&1{,}31102877714605987\\
{}[3\varpi=]\;&3{,}93308633143817961
\end{aligned}\]
\[\renewcommand{\arraystretch}{1.18}\begin{array}{r}
983\,271\,582\,859\,544\,902\\
15\,732\,345\,325\,752\,718\\
983\,271\,582\,859\,545\\ \hline
999\,987\,199\,768\,157\,165\\
9\,999\,871\,997\,682\\
2\,999\,961\,599\,305\\ \hline
1{,}000\,000\,199\,601\,754\,152
\end{array}\]
\clearpage
\[\renewcommand{\arraystretch}{1.18}\begin{array}{@{}r@{\;}r@{\qquad}r@{}}
[2\log_{10}e=0{,}]&86\,858\,896\,380\,650\,365&\\
&173\,717\,792\,761\,301&\\ \cline{2-2}
&86\,685\,178\,587\,889\,064&\\
&760\,015\,343\,33\phantom{0}&\\
&1\,737\,177\,927&\\
&65\,144\phantom{000}&\\
&869\phantom{000}&866\,859\,404\,064\\ \cline{2-2}
&86\,685\,940\,406\,423\phantom{000}&1\,733\,718\,808\\ \cline{3-3}
&8\,651\,333\phantom{000}&865\,125\,685\,256\\ \cline{2-2}
&[86\,685\,931\,755\,090]\phantom{000}&758\phantom{000}\\ \cline{3-3}
&&[865\,133\,26]\phantom{000}
\end{array}\]
\clearpage\SourceLeaf{343}\ArithmeticHead{336}{L}{ANALYSIS.\quad NACHLASS.}
\[\renewcommand{\arraystretch}{1.18}\begin{array}{@{}r@{\;}rrrr@{}}
339\;{.}\;{.}\;{.}&2{,}53019&96982&03082&16009\\
{}[225]\;{.}\;{.}\;{.}&2{,}35218&25181&11362&48416\\ \cline{2-5}
&4{,}88238&22163&14444&64425\\
&0&56457&91567&17660\\ \cline{2-5}
&88238&78621&06011&82085\\
&&866&85931&75509\\ \cline{2-5}
\log\varpi=&0{,}11761&22245&79919&93424\\
\multicolumn{4}{@{}l}{\text{eine scharfe Rechnung gab}}&94762
\end{array}\]
\begin{center}\rule{.22\linewidth}{.4pt}\end{center}
\[\renewcommand{\arraystretch}{1.2}\begin{array}{@{}r@{\quad}rrrrrr@{}}
\varpi\;{.}\;{.}\;{.}&0{,}11761&22245&79919&94586&18858&564\\
\surd\pi\;{.}\;{.}\;{.}&0{,}24857&49363&47066&92717&56368&\\[1em]
\pi^{\frac14}\varpi^{\frac12}&0{,}18309&35804&63493&43651&87610&\\
\overset{\scriptscriptstyle4}{\surd}8&0{,}22577&24967&47985&89641&03042&\\ \cline{2-6}
\Pi\,0{,}25&9{,}95732&10837&15507&54010&84568&\\[.5em]
\left[\frac{\surd\pi}{\varpi}\right]&0{,}13096&27117&67146&98131&37517&\\[.5em] \cline{2-6}
\Pi-0{,}25&0{,}08828&37954&82654&52142&22085&\\
0{,}75&9{,}87506&12633&91700&04686&75501&\\ \cline{2-6}
\Pi\,0{,}75&9{,}96334&50588&74354&56828&97586&
\end{array}\]
\[ {.}\;{.}\;{.}\;{.}\;{.}\;{.}\;{.}\;{.}\;{.}\;{.}\;{.}\;{.}\;{.}\;{.}\;{.}\;{.}\;{.}\;{.} \]
\noindent [Es folgt durch wiederholte Anwendung desselben Verfahrens die Berechnung der Logarithmen von $\Pi\,1{,}75$, $\Pi\,2{,}75$, $ {.}\;{.}\;{.}\;\Pi\,11{,}75$ auf je 20 Dezimalstellen. Dann von $\Pi\frac14$ ausgehend die Berechnung der Logarithmen von $\Pi\,1{,}25$, $\Pi\,2{,}25$, $ {.}\;{.}\;{.}\;\Pi\,12{,}25$.]
\par\bigskip\noindent{\small[S. 247]}\par
\[\Pi\,7{,}5=\frac{1001{.}8100}{1024}\,\surd\pi\]
\[\renewcommand{\arraystretch}{1.25}\begin{array}{@{}r@{\quad}rrrr@{}}
\frac{1001}{1024}&9{,}99013&41208&39506&68853\\[.5em]
{}[8100]&3{,}90848&50188&78649&74918\\ \cline{2-5}
&3{,}89861&91397&18156&43771\\
{}[\surd\pi]&0{,}24857&49363&47066&92718\\ \cline{2-5}
\Pi\,7{,}5&4{,}14719&40760&65223&36489
\end{array}\]
\clearpage\SourceLeaf{344}\ArithmeticHead{337}{R}{{\scriptsize ZUR THEORIE DER REIHE $F(\alpha,\beta,\gamma,x)$.\quad II.}}
\noindent [Es folgt die Berechnung der Logarithmen von $\Pi\,8{,}5$, $\Pi\,9{,}5$, $\Pi\,10{,}5$, $\Pi\,11{,}5$ dann von $\Pi\,8{,}25$ bis zu $\Pi\,13{,}00$, fortschreitend je um $0{,}25$, alles auf 20 Dezimalstellen.]
\endgroup
